\section{Introdução}\label{sec:introducao}

Este relatório tem como objetivo documentar o desenvolvimento de uma aplicação de supervisão e controle utilizando o software \textbf{CODESYS v3.5}\cite{codesysV35SP21} e o supervisório \textbf{InduSoft Web Studio}. O projeto contempla a automação de uma planta de produção de Dry Martini, abrangendo desde a dosagem e mistura dos ingredientes (Área 01) até o envase e transporte do produto final (Área 02).

O sistema desenvolvido integra lógicas sequenciais (SFC), controle contínuo (Ladder), algoritmos matemáticos (Texto Estruturado) e comunicação via protocolo \textbf{OPC UA}. A aplicação foi estruturada em tarefas multitarefa para garantir a fidelidade da simulação física e a precisão do controle.

Este documento detalha a divisão das tarefas, a lógica implementada em cada Unidade de Organização de Programa (POU) e a interface desenvolvida para operação.

Todo o projeto pode ser encontrado no Git-Hub atráves do url:  \url{https://github.com/Rian-Rero/Dry-Martini}\cite{dryMartiniRepo}